% LaTeX preamble for the paper PDF build (step 47; see the Makefile's step_47 target).
%
% The manuscript is written in Markdown with real Unicode mathematical and typographic characters
% (>=, <=, a check mark, superscripts, Greek). Tectonic renders through XeTeX with Latin Modern,
% which has no glyph for several of them and silently DROPS them -- a "roc_auc >= 0.75" in the
% source would reach the PDF as "roc_auc  0.75", turning a threshold into a typo without any error.
%
% Mapping them here rather than switching mainfont to a wide-coverage face (DejaVu) keeps the
% manuscript in Latin Modern, which is what the text is set in.
\usepackage{newunicodechar}
\usepackage{amssymb}
\usepackage{textcomp}

\newunicodechar{≥}{\ensuremath{\geq}}
\newunicodechar{≤}{\ensuremath{\leq}}
\newunicodechar{−}{\ensuremath{-}}
\newunicodechar{×}{\ensuremath{\times}}
\newunicodechar{Δ}{\ensuremath{\Delta}}
\newunicodechar{μ}{\ensuremath{\mu}}
\newunicodechar{✓}{\ensuremath{\checkmark}}
\newunicodechar{⁻}{\textsuperscript{-}}
\newunicodechar{⁷}{\textsuperscript{7}}
\newunicodechar{¹}{\textsuperscript{1}}
\newunicodechar{²}{\textsuperscript{2}}

% Table 1 is seven columns of prose and overruns the text block badly at full size; the benchmark
% table is wider still. Shrinking table type is what keeps them on the page without rewriting the
% manuscript's content to fit its own typesetting.
\let\oldlongtable\longtable
\def\longtable{\footnotesize\oldlongtable}

% Constrain every included figure to the text width HERE rather than with a pandoc `{width=95%}`
% attribute on each image tag. The manuscript is also read as Markdown on the repository page,
% where that attribute is not image syntax and renders as literal "{width=95%}" text after every
% figure. Setting it in the preamble keeps paper/README.md portable and the PDF correctly sized.
% The figures are 300-DPI Plotly exports ~3900 px wide, so without this they overrun the page.
%
% graphicx is loaded explicitly because pandoc only emits it for a document that actually contains
% an image. The supplement contains none, so \setkeys{Gin} alone fails there with the decidedly
% unhelpful "Package keyval Error: width undefined" -- this header is shared by both documents.
\usepackage{graphicx}
\setkeys{Gin}{width=0.95\linewidth,keepaspectratio}

% Suppress LaTeX's own "Figure N:" counter on captions. The manuscript's figure numbers come from
% the generator's filenames (figures/NN_*.png) and are what the prose cites -- "(Figure 14)" -- so
% each caption carries its number explicitly as text. LaTeX's counter numbers by POSITION instead,
% which does not agree: the figures appear in citation order, and citation order is not numeric
% order here (Figure 14 is cited in Section 3.2, before Figure 3 in Section 3.3). Left on, every
% caption reads "Figure 3: Figure 14. ..." -- numbered twice, and disagreeing with itself.
\usepackage{caption}
\captionsetup[figure]{labelformat=empty}

% Figures land where they are cited rather than drifting pages away to a float page: the
% manuscript's Results sections read as prose-then-evidence, and LaTeX's default float placement
% breaks that pairing in a document whose sections are single dense paragraphs.
\usepackage{float}
\floatplacement{figure}{H}
